\section{Analytic reconstruction for nonlinear observability systems}\label{sec:abstract-construction}

The purpose of this section is to prove \cref{thmabstractanalyticintro}. For the reader's convenience, we~will briefly outline its proof, following the Galerkin decomposition introduced in Hale-Raugel \cite{HR03}. The key replacement will be the observability estimate that allows to reconstruct the state from the observation, at least for the high-frequency part.

Let $T>0$ and fix $\sigma\in [0, 1]$. From \cref{assumAA}, we~introduce the low and high-frequency projections $\P_n=\mathbbm{1}_{[0, n]}\big((AA^*)^{1/2}\big)$ and $\Q_n=I-\P_n$, respectively. Let $U=U(t)$ be a mild solution of \eqref{UCabstractT*intro} in $C^0([0, T], X^\sigma)$ and suppose $H_1=0$, $H_2=0$ for simplicity. Let us consider the splitting
\begin{align*}
U(t)=\P_n U(t)+\Q_nU(t)=V(t)+W(t),
\end{align*}
where $\big(V(t), W(t)\big)$ solves the following system
\begin{align*}
\left\{\begin{array}{rl}
\partial_t V(t)&=AV(t)+\P_n F(V+W), \\[3pt]
\partial_t W(t)&=AW(t)+\Q_n F(V+W), \\[3pt]
\bC V(t)&=-\bC W(t).
\end{array}\right.
\end{align*}
By Duhamel's formula, the high-frequency component $W$ can be written as
\begin{align*}
W(t)=e^{tA}W(0)+\int_0^t e^{(t-s)A}\Q_n F\big(V(s)+W(s)\big)ds.
\end{align*}
The observation condition $\bC V=-\bC W$ suggests that given $V$, we~can reconstruct $W$ by considering the corresponding nonlinear observability system. Indeed, according to \cref{assumCC}, the observability of the linear semigroup $t\in [0, T]\mto e^{tA}$ enables us to construct an initial condition $W(0)$ solely in terms of an observation. Forgetting first about the source term given by the nonlinearity, that would allow to reconstruct $W(0)$ in terms of the observation of $W$, which is $\bC W=-\bC V$. \cref{lmCauchyobs} below provides a generalization of this reconstruction problem when source terms are present. In~this context, the first part of \cref{assumFholom}, namely, that $F$ is a nonlinear map from bounded sets of $X^\sigma$ into $X^{\sigma+\veps}$, will allow, at frequency sufficiently large, to consider the nonlinearity as a perturbation and to complete this reconstruction procedure. More precisely, this will imply the existence of a nonlinear map $\mc{N}$ such that $W(0)=\mc{N}(V)$. Consequently, we~have the formula $W=\Phi_V(W)$, where $\Phi_V: C^0([0, T], \Q_nX^\sigma)\to C^0([0, T], \Q_nX^\sigma)$ is given by
\begin{align*}
\Phi_V(W)(t)=e^{tA}\mc{N}V(\cdot)+\int_0^t e^{(t-s)A}\Q_n F\big(V(s)+W(s)\big)ds,\ t\in [0, T].
\end{align*}
This suggests that we can find the high-frequency component $W^*$ as a fixed point with the low-frequency component $V$ as an input.

At this point, the solution $U$ can be represented as $U(t)=V(t)+W^*(V)(t)$, where~$V$ solves
\begin{align}\label{intro:eq-low-split}
\partial_t V(t)=AV+\P_n F(V+W^*(V)).
\end{align}
To demonstrate that $t\in (0, T)\mto U(t)\in X^\sigma$ is analytic, the second part of \cref{assumFholom}, namely, that $F$ admits a holomorphic extension, is essential. This will be achieved by establishing that $t\mto V(t)$ and $V\mto W^*(V)$ are both analytic maps. If instead, we~consider \eqref{intro:eq-low-split} as a differential equation on the space Banach space $C^0([0, T], \P_n X^\sigma)$, classical ODEs theory imply that $t\mto V(t)$ is as smooth as $F$, and therefore analytic. The uniform contraction principle further ensures that $W^*$ depends analytically on $V$, from which the result follows.

In what follows, we~develop these ideas towards the proof of the main theorem. Throughout the section, we~will prove several intermediate results under weaker assumptions, some of which are of independent interest. For instance, we~mention that \cref{prop:finite-det-modes} (Finite determining modes) and \cref{prop:abs-prop-reg} (Propagation of regularity) below do not require analyticity. Nevertheless, the hurried reader could consider all the following results under Assumptions~\ref{assumAA}, \ref{assumCC}, \ref{assumFholom} and \ref{assumcommu}, which will always contain the intermediate assumptions made throughout this section.

\Subsection{Linear reconstruction, determining modes and propagation of regularity} In this section, we~will revisit the assumptions outlined in the introduction, organizing them according to the different results we aim to establish in order to prove our main result.

From now on, we~work under \cref{assumAA}. We~will now list some consequences of such an assumption. That $A^*A=-A^2$ is non-negative self-adjoint, allows us to define the Hilbert space $X^{\sigma}=D((A^*A)^{\sigma/2})$ for any $\sigma\in\R$. Note that the assumptions imply $X^{\sigma+\veps}\hookrightarrow X^{\sigma} $ for any $\veps>0$. Unless specifically noticed, we~will often omit the embedding $\iota: X^{\sigma+\veps}\to X^{\sigma}$.

By the spectral theorem, and since $A^*A$ has a compact resolvent, thus, the spectrum of $A^*A$ is real and discrete, allowing us to construct an orthonormal basis of eigenvectors of $A^*A$ in $X$, denoted by $(E_j)_{j\in\N}$ and associated to the nonnegative eigenvalues $(\lambda_j)_{j\in \N}$ (ranged increasingly) with $\lambda_j \to+\infty$ as $j\to +\infty$. We~introduce the high-frequency projectors $\Q_n$ on the space $\overline{\text{Span}\{E_j\}_{j\geq n}}$ and then we set the low-frequency projection $\P_n=I-\Q_n$. Note that $A$ commutes with $\P_n$ and $A\P_n$ is a bounded operator of $X^{\sigma}$ to itself with norm $\inn{\ld_n}$.

The parameter $\sigma$ will be fixed from now on. We~will use the notation $\P_n X^{\sigma}$ or $\Q_n X^{\sigma}$ for that related image of the Hilbert space endowed with the topology of~$X^{\sigma}$. The restriction of the embedding $\iota$ to $\Q_n X^{\sigma+\veps}$(denoted with the same name) $\iota: \Q_n X^{\sigma+\veps}\to \Q_nX^{\sigma}$ has norm $\inn{\lambda_n}^{-\veps}$.

By the spectral theorem, since the spectrum of $A$ is purely imaginary, we~can define $e^{tA}$ for any $t\in \R$, together with the estimate $\norm{e^{tA}U_0}_{X^{\sigma}}=\norm{U_0}_{X^{\sigma}}$. Also, $e^{tA}$~commutes with $\P_n$ and $\Q_n$.

\subsubsection{Linear reconstruction} We will now introduce the following assumption regarding the observability of the semigroup generated by $A$, which corresponds to the first part of Assumption~\ref{assumCC}.
\begin{assump}{2a}
\label{assumC}Let $\bC\in\mc{L}(X^\sigma, X^\sigma)$ be an observation operator. We~assume that $t\mto e^{tA}$ is observable on $[0,T]$, namely, there exists a constant $\mathfrak{C}_{obs}>0$ such that
\begin{align}\label{observabstract}
\norm{W_0}_{X^\sigma}^2\leq \mathfrak{C}_{\textup{obs}}^2\int_0^T \norm{\bC e^{tA}W_0}_{X^\sigma}^2dt,\quad \forall W_0\in X^\sigma.
\end{align}
\end{assump}
Let $\mc{O}\in \mc{L}(X^\sigma, L^2([0, T], X^\sigma))$, defined by $\mc{O}:=\bC e^{\cdot A}$, be the observation operator of linear waves and $\mc{O}_n:=\mc{O}|_{\Q_nX^\sigma}$ the \emph{high-frequency} observation operator. The observability inequality \eqref{observabstract} can be written
\begin{align}\label{observabstractO}
\norm{W_0}_{X^\sigma}&\leq \mathfrak{C}_{\textup{obs}}\nor{\O W_{0}}{L^{2}([0,T],X^\sigma)},\quad\;\; \forall W_0\in X^\sigma,\\
\label{observabstractOn} \norm{W_0}_{X^\sigma}&\leq \mathfrak{C}_{\textup{obs}}\nor{\O_{n} W_{0}}{L^{2}([0,T],X^\sigma)},\quad \forall W_0\in \Q_n X^\sigma,
\end{align}
where the last inequality is uniform in $n\in \N$.
It implies that $\mc{O}$ is injective and it has a closed range. Moreover, since $\Q_nX^\sigma$ is closed in $X^\sigma$, then the \emph{high-frequency} observation operator $\mc{O}_n:=\mc{O}|_{\Q_nX^\sigma}$ has closed range as well, which allows us to define $\Pi_n$ as the orthogonal (according to the natural scalar product in $L^2([0, T], X^\sigma)$) projection onto its image $\text{Im}(\mc{O}_n)\subset L^2([0, T], X^\sigma)$. From now on, we~equip $\Y_n:=\text{Im}(\O_n)$ with the induced topology from $L^2([0, T], X^\sigma)$ which makes it a Banach space. By~\eqref{observabstractOn}, we~know that $\O_n: \Q_n X^\sigma\to \Y_n$ is a bijection, and hence, $\Y_n$ is closed, a bounded reconstruction operator $\O_n^{-1}: \Y_n\to \Q_n X^\sigma$ exists.

By applying the observability inequality \eqref{observabstractOn}, consequence of \cref{assumC}, we~get for any $y\in \Y_n \subset L^2([0, T], X^\sigma)$,
\begin{align}
\label{boundOn-1}
\norm{\O_n^{-1}y}_{X^\sigma}&\leq \mathfrak{C}_{\textup{obs}}\norm{\O_{n}\O_n^{-1}y}_{L^2([0, T], X^\sigma)}=\mathfrak{C}_{\textup{obs}}\norm{y}_{L^2([0, T], X^\sigma)},
\end{align}
where again, the inequality is uniform in $n\in \N$.

To ease notation, we~consider the operator $\mc{I}(t): g\mto \int_0^t e^{(t-s)A}g(s)ds$ and we denote $\mc{I}(\cdot)$ when the operator is seen with value in $C^{0}([0,T],Y)$ for a suitable Banach space. The above construction will enable us to solve an \emph{observability Cauchy problem}, which is the content of the following lemma.

\begin{lemma}
\label{lmCauchyobs}
There exists $C(T,\mathfrak{C}_{\textup{obs}},\nor{\bC}{\mc{L}(X^\sigma)})>0$ such that for any $n\in\N$, $H\in L^{1}([0,T],X^{\sigma})$ and $G\in L^{2}([0,T],X^{\sigma})$, there exists a unique $W\in C^{0}([0,T],\Q_{n}X^{\sigma})$ solution of
\begin{align}\label{solprojectlin}
\left\{\begin{array}{rl}
\partial_t W(t)&=AW(t)+\Q_{n}H, \\[3pt]
\Pi_n\bC W &=\Pi_{n}G.
\end{array}\right.
\end{align}
It satisfies $W(0)= W_0:=\O_n^{-1}\Pi_n\left[G-\bC\mc{I}(\cdot)\Q_n H\right]$ and is given by $W(t)=e^{tA}W_{0}+\mc{I}(t)\Q_nH$. We~denote $\FL$ the associated linear operator defined by $\FL(G,H):=W$. Moreover, we~have the estimate
\begin{align}
\label{estimFL}
\nor{\FL(G,H)}{C^{0}([0,T],\Q_{n}X^{\sigma})}\leq C \nor{\Pi_{n}G}{L^{2}([0,T],X^{\sigma})}+C\nor{\Q_{n} H}{L^{1}([0,T],X^{\sigma})}
\end{align}
\end{lemma}
\begin{proof}
To be solution of the first line of \eqref{solprojectlin}, it~is equivalent to be written as a Duhamel formula $W(t)=e^{tA}W_{0}+\mc{I}(t)\Q_nH$ for some $W_{0}\in \Q_{n}X^{\sigma}$. So, we~only need to compute $W_{0}$. With the previous formula, we~have $W\in C^{0}([0,T],\Q_{n}X^{\sigma})\subset L^{2}([0,T],\Q_{n}X^{\sigma})$ and we can compute
\begin{align*}
\Pi_n\bC W= \Pi_n\bC \left[e^{\cdot A}W_{0}+\mc{I}(\cdot)\Q_nH \right]=\Pi_n\bC \left[e^{\cdot A}W_{0} \right]+\Pi_n\bC \left[\mc{I}(\cdot)\Q_nH \right].
\end{align*}
Note that if $W_{0}\in \Q_{n}X^{\sigma}$, then $\bC\left[e^{\cdot A}W_{0}\right]=\O W_{0}=\O_{n} W_{0}$ and therefore, by~definition of $\Pi_{n}$,
\begin{align*}
\Pi_{n}\bC\left[e^{\cdot A}W_{0}\right]=\Pi_{n}\O_{n} W_{0}=\O_{n} W_{0}.
\end{align*}

In particular, since both belong to $\mathcal{Y}_n$, we~want $ \Pi_n\bC W=\Pi_{n}G$, we~should have
\begin{align*}
\O_n^{-1}\Pi_{n}G= \O_n^{-1} \Pi_n\bC W &=\O_n^{-1}\O_{n} W_{0}+\O_n^{-1}\Pi_n\bC \left[\mc{I}(\cdot)\Q_nH \right]\\ &=W_{0}+\O_n^{-1}\Pi_n\bC[\mc{I}(\cdot)\Q_nH].
\end{align*}
This gives that $W_{0}$ should be defined by the formula given in the lemma. It~indeed belongs to $\Q_{n}X^{\sigma}$ and therefore $W$, as defined, satisfies the second line of \eqref{solprojectlin} by reproducing the same computation backward, that is\vspace*{-3pt}
\begin{align*}
\Pi_n\bC W&=\Pi_n\bC\left[e^{\cdot A}W_{0}+\mc{I}(\cdot)\Q_nH\right] =\Pi_{n} \O_{n}W_{0}+\Pi_n\bC\mc{I}(\cdot)\Q_nH\\
&=\Pi_{n} \O_{n}\O_n^{-1}\Pi_n\left[G-\bC\mc{I}(\cdot)\Q_n H\right]+\Pi_n\bC\mc{I}(\cdot)\Q_nH\\
&=\Pi_n\left[G-\bC\mc{I}(\cdot)\Q_n H\right]+\Pi_n\bC\mc{I}(\cdot)\Q_nH=\Pi_{n}G.
\end{align*}
The uniqueness could actually be obtained from the unique definition of $W_{0}$ that we obtained, but we prefer to give a precise proof. We~consider the difference $R=W_{1}-W_{2}\in C^{0}([0,T],\Q_{n}X^{\sigma})$ between two such solutions $W_{1}$ and $W_{2}$. It~satisfies\vspace*{-3pt}
\begin{align*}
\left\{\begin{array}{rl}
\partial_t R(t)&=AR(t), \\[3pt]
\Pi_n\bC R &=0.
\end{array}\right.
\end{align*}
That is $R(t)=e^{tA}R(0)$ and $\Pi_n\bC R=\Pi_{n}\O R(0)=\Pi_{n}\O_{n}R(0)=\O_{n}R(0)$. In~particular, $R(0)=0$ by injectivity of $\O_{n}$.

Concerning the estimates, since $A$ is skew-adjoint on $X^{\sigma}$, standard semigroup estimates give\vspace*{-3pt}
\begin{align}
\nor{W}{C^{0}([0,T],\Q_{n}X^{\sigma})}\leq \nor{W_{0}}{X^{\sigma}}+\nor{\Q_{n} H}{L^{1}([0,T],X^{\sigma})}.
\end{align}
So, we~need to estimate $W_{0}$. For any $\widetilde{G}\in L^2([0, T], X^\sigma)$, applying \eqref{boundOn-1} to $\Pi_n \widetilde{G}\in \Y_n$, we~have\vspace*{-3pt}
\begin{align*}
\norm{\O_n^{-1}\Pi_n \widetilde{G}}_{X^\sigma}&\leq \mathfrak{C}_{\textup{obs}}\norm{\Pi_n\widetilde{G}}_{L^2([0, T], X^\sigma)}.
\end{align*}
In particular, we~can estimate $W_{0}$ by\vspace*{-3pt}
\begin{align*}
\nor{W_{0}}{X^{\sigma}}&= \nor{\O_n^{-1}\Pi_n\left[G-\bC\mc{I}(\cdot)\Q_n H\right]}{X^{\sigma}}\\ &\leq \mathfrak{C}_{\textup{obs}}\norm{\Pi_nG}_{L^2([0, T], X^\sigma)}+\mathfrak{C}_{\textup{obs}}\norm{\Pi_n\bC\mc{I}(\cdot)\Q_n H}_{L^2([0, T], X^\sigma)}.
\end{align*}
We can finally estimate by the unitarity of $\Pi_{n}$ and Hölder inequality in time\vspace*{-3pt}
\begin{align*}
\norm{\Pi_n\bC\mc{I}(\cdot)\Q_n H}_{L^2([0, T], X^\sigma)}&\leq \norm{\bC\mc{I}(\cdot)\Q_n H}_{L^2([0, T], X^\sigma)}\\ &\leq \nor{\bC}{\mc{L}(X^\sigma)}\norm{\mc{I}(\cdot)\Q_n H}_{L^2([0, T], X^\sigma)}\\
&\leq T^{1/2}\nor{\bC}{\mc{L}(X^\sigma)}\norm{\mc{I}(\cdot)\Q_n H}_{L^{\infty}([0, T], X^\sigma)}\\ &\leq T^{1/2}\nor{\bC}{\mc{L}(X^\sigma)}\norm{\Q_n H}_{L^{1}([0, T], X^\sigma)}.
\end{align*}
Recollecting the previous estimates, we~have finally proved\vspace*{-8pt}
\begin{multline*}
\nor{W}{C^{0}([0,T],\Q_{n}X^{\sigma})}\\\leq \mathfrak{C}_{\textup{obs}} \nor{\Pi_{n}G}{L^{2}([0,T],X^{\sigma})}+\bigl(1+T^{1/2}\mathfrak{C}_{\textup{obs}}\nor{\bC}{\mc{L}(X^\sigma)}\bigr)\nor{\Q_{n} H}{L^{1}([0,T],X^{\sigma})}.\qedhere
\end{multline*}
\end{proof}
\begin{remark}
It will be very important for what follows that the constant $C$ involved in the previous lemma is independent of $n\in\N$.
\end{remark}

\subsubsection{Finite determining modes}
As a first direct consequence of the previous result, we~can get a finite determining mode result: two solutions of a nonlinear equation with the same observation and the same low frequency are the same. This result will not be used directly later, but can be considered as an easier version of what will follow where we will actually construct the reconstruction operator and study its regularity.

We make the following assumption on the nonlinearity $F$, akin to a compactness property. This corresponds to the first part of Assumption~\ref{assumFholom}.

\begin{assump}{3a}
\label{assumF}
$F$ is a nonlinear operator from $\mathbb{B}_{4R_{0}}(X^{\sigma})$ to $X^{\sigma+\veps}$ for some $R_{0}>0$, $\sigma>0$ and $\veps>0$. Moreover, there exists $C>0$ such that
\begin{align}
\label{boundF}
\nor{F(U_{0})}{X^{\sigma+\veps}}\leq C,\quad \nor{F(U_{0})-F(V_{0})}{X^{\sigma+\veps}}\leq C \nor{U_{0}-V_{0}}{X^{\sigma}}
\end{align}
for any $U_{0},~V_{0}\in \mathbb{B}_{4R_{0}}(X^{\sigma})$.
\end{assump}

Note that the second bound implies the first one with another constant. We~show the following property of finite determining modes.

\begin{proposition}\label{prop:finite-det-modes}
Let $R_{0}>0$. Under Assumptions \ref{assumAA}, \ref{assumC} and \ref{assumF} (with $R_{0}$), there exists $n\in \N$ such that the following holds. Let $H\in L^1([0,T],X^{\sigma})$ and $G\in L^2([0, T], X^\sigma)$.

Let $U(t)$ and $\widetilde{U}(t)$ be two solutions on $(0,T)$ of
\begin{align*}
\left\{\begin{aligned}
\partial_t U&=AU+F(U)+H, &&\text{on } (0, T), \\
\bC U(t)&=G(t), && \text{for } t\in (0, T),
\end{aligned}\right.
\end{align*}
such that $\norm{U(t)}_{X^\sigma}\leq R_{0}$ and $\norm{\widetilde{U}(t)}_{X^\sigma}\leq R_{0}$ for all $t\in [0, T]$. If $\P_n U(t)=\P_n \widetilde{U}(t)$ for all times $t\in [0, T]$, then $U(t)\equiv \widetilde{U}(t)$ for all $t\in [0, T]$.
\end{proposition}
\begin{proof}
By assumption, $\P_{n}U=\P_{n}\widetilde{U}$ as maps in $\B_{R_{0}}^{[0,T]}(X^{\sigma})$. Let us consider the difference of solutions $Z=U-\widetilde{U}$. It~satisfies
\begin{align*}
\left\{\begin{aligned}
\partial_t Z&=AZ+F(U)-F(\widetilde{U}), \\
\bC Z&=0.
\end{aligned}\right.
\end{align*}
Moreover, we~have $\P_n Z=0$, that is $Z\in C^{0}([0,T],\Q_{n}X^{\sigma})$ and therefore, applying~$\Q_n$, it~also satisfies
\begin{align*}
\left\{\begin{aligned}
\partial_t Z&=AZ+\Q_n \bigl(F(U)-F(\widetilde{U})\bigr), \\
\Pi_{n}\bC Z&=0.
\end{aligned}\right.
\end{align*}
In particular, we~are in the situation of \cref{lmCauchyobs} and $Z=\FL(0,\Q_n \big(F(U)-F(\widetilde{U})\big))$, so that estimate \eqref{estimFL} gives
\begin{align*}
\nor{Z}{C^{0}([0,T],\Q_{n}X^{\sigma})}&\leq C \bnor{\Q_n \bigl(F(U)-F(\widetilde{U})\bigr)}{L^{1}([0,T],X^{\sigma})}\\
&\leq \dfrac{C}{\inn{\ld_n}^\veps} \bnor{F(U)-F(\widetilde{U})}{L^{1}([0,T],X^{\sigma+\veps})}.
\end{align*}

By hypothesis, $U$, $\widetilde{U}\in \B_{R_0}^{[0, T]}(X^\sigma)$, and \cref{assumF} implies that $F(U)$, $F(\widetilde{U})\in \B_{C}^{[0,T]}(X^{\sigma+\veps})$ with the following Lipschitz estimate
\begin{align*}
\nor{Z}{C^{0}([0,T],X^{\sigma})}\leq \dfrac{C}{\inn{\ld_n}^\veps}\, \bnor{U-\widetilde{U}}{L^{1}([0,T],X^{\sigma})}\leq \dfrac{CT}{\inn{\ld_n}^\veps} \nor{Z}{C^{0}([0,T],X^{\sigma})}.
\end{align*}
Adjusting $n$ if necessary so that $\sfrac{CT}{\inn{\ld_n}^\veps}<1$, we~conclude that $Z=0$, which is $U= \widetilde{U}$.
\end{proof}

\subsubsection{Propagation of regularity} We now turn our attention to the regularity of the nonlinear system, for which we prove a propagation-of-regularity result under suitable assumptions. This result will later be related to a sort of uniformity for the splitting in the Galerkin procedure.

Building upon \cref{assumC}, we~make the following assumption, which will allow us to control the semigroup generated by $A$ in a slightly more regular space:
\begin{align}\label{observabstractveps}
\norm{W_0}_{X^{\sigma+\veps}}^2\leq \mathfrak{C}_{\textup{obs}}^2\int_0^T \norm{\bC e^{tA}W_0}_{X^{\sigma+\veps}}^2dt,\quad \forall W_0\in X^{\sigma+\veps}.
\end{align}
In that context, \cref{assumCC} is the conjunction of \cref{assumC} and \eqref{observabstractveps}, that is, the observability at both levels or regularity $X^{\sigma}$ and $X^{\sigma+\veps}$.

\begin{proposition}\label{prop:abs-prop-reg}
Let $R_{0}>0$, $R_{1}>0$ and $A$, $\bC$ and $F$ satisfying Assumptions~\ref{assumAA},~\ref{assumCC} and~\ref{assumF}, respectively. Moreover, assume that the pair $(A, \bC)$ satisfies \cref{assumcommu}. Then there exists $R_{2}>0$ such that for any
$U\in \B_{R_{0}}^{[0,T]}(X^{\sigma})$, $H_{1}\in \B_{R_{0}}^{[0,T]}(X^{\sigma})$ and $H_{2}\in \B_{R_{1}}^{[0,T]}(X^{\sigma+\veps})$ that satisfy \eqref{UCabstractT*intro} on $[0,T]$, we~have $U\in \B_{R_{2}}^{[0,T]}(X^{\sigma+\veps})$.
\end{proposition}
\begin{proof}

We know that $(AA^*)^{1/2}: X^1\to X$ admits a unique restriction such that $(AA^*)^{1/2}: X^{1+\sigma}\to X^{\sigma}$ is a linear continuous operator. Let us call such extension~$\A_\sigma$. Furthermore:
\begin{itemize}
\item From \cref{assumAA}, $\A_\sigma$ has compact resolvent.
\item $\A_\sigma$ and $(AA^*)^{1/2}$ have the same spectrum, hence the same resolvent set.
\end{itemize}
For simplicity, we~keep the same notation $\A$ for the same operator acting on different spaces. Since $\A$ is non-negative, its resolvent set contains $\R_-$ and for $n\in \N^*$, for any $s>0$, we~have a well-defined smoothing operator $\mc{J}_n=(I+\frac{1}{n}\A^s)^{-1}\in \mc{L}(X^\sigma, X^{\sigma+s})$ with the uniform bound $\norm{\mc{J}_n}_{\mc{L}(X^\sigma, X^{\sigma+s})}\leq n$. We~also have the uniform bound $\norm{\mc{J}_n}_{\mc{L}(X^\sigma)}\leq 1$ and the same estimate holds in $\mc{L}(X^{\sigma+\veps})$. Following \cite[Prop.\,2.3.4]{TW:09}, we~can see that $\J_n \phi\underset{n\to+\infty}{\to} \phi$ for any $\phi\in X^\sigma$.

Let us consider $U^n(t)=\J_nU(t)$ and observe that $U^n$ is uniformly bounded in $X^\sigma$ by some constant $C>0$, which is independent of $n\in \N$. By~Duhamel's formula, let~us split the solution into its linear and nonlinear part as follows
\begin{align*}
U^n(t)&\hphantom{:}=e^{tA}U_0^n+\int_0^t e^{(t-s)A}\big(\J_nF(U(s)+H_1(s))+\J_nH_2(s)\big)ds\\ &:=U_{\lin}^n(t)+U_{\Nlin}^n(t),
\end{align*}
where $U_0^n:=\J_n U_0$. Observe that we have used that $e^{tA}$ and $\J_n$ commute. Since \hbox{$H_2\in \B_{R_{1}}^{[0, T]}(X^{\sigma+\veps})$} and $F$ satisfies \cref{assumF}, we~readily get that $U_{\Nlin}^n$ is uniformly bounded in $L^\infty([0, T], X^{\sigma+\veps})$ by some constants depending on $R_{1}$ and the constant $C$ in \cref{assumF}. The latter property and \cref{assumCC}, allows us to treat the linear part by employing the observability inequality followed by the triangle inequality
\begin{align*}
\norm{U_{\lin}^n(t)}_{X^{\sigma+\veps}}^2&\leq \norm{e^{tA}}_{\mc{L}(X^{\sigma+\veps})}^2 \norm{U_0^n}_{X^{\sigma+\veps}}^2\\
&\leq \mathfrak{C}_{\textup{obs}}^2\int_0^T \norm{\bC U_{\lin}^n(t)}_{X^{\sigma+\veps}}^2dt\\
&\leq 2\mathfrak{C}_{\textup{obs}}^2\int_0^T \norm{\bC U^n(t)}_{X^{\sigma+\veps}}^2dt+2\mathfrak{C}_{\textup{obs}}^2\int_0^T \norm{\bC U_{\Nlin}^n(t)}_{X^{\sigma+\veps}}^2dt.
\end{align*}
Since $\bC U\equiv 0$ on $[0, T]\times \M$, we~get that
\[
\bC U^n=\bC\mc{J}_n U=\mc{J}_n\bC U+[\mc{J}_n,\bC]U=[\mc{J}_n,\bC]U.
\]
Moreover, we~have $[\mc{J}_n,\bC]=\frac{1}{n}\mc{J}_n[\A^s,\bC]\mc{J}_n$. We~have seen that $\norm{\mc{J}_n}_{\mc{L}(X^\sigma)}\leq 1$ and $\norm{\frac{1}{n}\mc{J}_n}_{\mc{L}(X^{\sigma}, X^{\sigma+s})}\leq 1$, so we get, uniformly in $n$,
\begin{align*}
\norm{\bC U^n}_{L^\infty([0, T], X^{\sigma+\veps})}&=\norm{\mc{J}_n[\A^s,\bC]\Psfrac{\mc{J}_n}{n} U}_{L^\infty([0, T], X^{\sigma+\veps})}\\
&\leq R_0\norm{[\A^s,\bC]}_{\mc{L}(X^{\sigma+s}, X^{\sigma+\veps})}.
\end{align*}
Therefore, in~view of \cref{assumcommu}, the term $\norm{[\A^s,\bC]}_{\mc{L}(X^{\sigma+s}, X^{\sigma+\veps})}$ is bounded and so $U_{\lin}^n$ is uniformly bounded in $L^\infty([0, T], X^{\sigma+\veps})$. It~follows that $U^n$ is uniformly bounded in $L^\infty([0, T], X^{\sigma+\veps})$ by some constant $C>0$. Moreover, due to the fact that $\norm{U-U^n}_{L^\infty([0, T], X^\sigma)}\to 0$ and by weak compactness in the Hilbert space $X^{\sigma+\veps}$, we~get that $U(t)\in X^{\sigma+\veps}$ and
\begin{align*}
\norm{U(t)}_{X^{\sigma+\veps}}\leq \liminf \norm{U^n(t)}_{X^{\sigma+\veps}}\leq C,
\end{align*}
for any $t\in [0, T]$, showing that $U$ is uniformly bounded in $L^\infty([0, T], X^{\sigma+\veps})$. We~obtain that $U\in C^0([0, T], X^{\sigma+\veps})$ using Duhamel formula now that we know that, for instance, $U(0)\in X^{\sigma+\veps}$ and $F(U+H_1)\in L^\infty([0, T], X^{\sigma+\veps})$.
\end{proof}
\begin{remark}\label{rk:abs-prop-reg}
Looking at the above proof, in~regards to the nonlinearity, we~only need to ensure that the assignment $t\mto \int_0^t e^{(t-s)A}F(Z(s))ds$ defines a bounded map in $L^\infty([0, T], X^{\sigma+\veps})$. For instance, this was achieved here through the sole hypothesis that $F(Z)$ is bounded in $L^\infty([0, T], X^{\sigma+\veps})$ for $Z$ bounded in $L^\infty([0, T], X^\sigma)$.
\end{remark}

\subsection{An abstract frequency-based reconstruction operator} The objective of this section is to prove that it is possible to reconstruct the high-frequency component for the solutions of our nonlinear system, with the low-frequency component as an input. Furthermore, this reconstruction can be made in a holomorphic way. For most of the facts about complex analysis in Banach spaces, our main reference is \cite{Muj86}.

Building upon \cref{assumF}, we~further assume that there exists $\delta>0$ such that~$F$ has a holomorphic extension from $\mathbb{B}_{4R_{0},2\delta}(X^{\sigma})$ to $X^{\sigma+\veps}_{\mathbb{C}}$ for some $R_{0}>0$, $\sigma\geq 0$ and $\delta>0$. Moreover, there exists $C>0$ such that
\begin{align}
\label{boundFanalytic}
\nor{F(U_{0})}{X^{\sigma+\veps}_{\mathbb{C}}}\leq C,\quad \nor{F(U_{0})-F(V_{0})}{X^{\sigma+\veps}_{\mathbb{C}}}\leq C \nor{U_{0}-V_{0}}{X^{\sigma}_{\mathbb{C}}}
\end{align}
for any $U_{0},~V_{0}\in \mathbb{B}_{4R_{0},2\delta}(X^{\sigma})$. In~this context, \cref{assumFholom} is the conjunction of \cref{assumF} and the previous assumption.

The main result of this section is the following.

\begin{theorem}
\label{thmabstract}Let $R_{0}$ and $R_{1}>0$. Under Assumptions \ref{assumAA}, \ref{assumCC}, \ref{assumF} (with $R_{0}$) and \ref{assumcommu}, there exist $n\in \N$ and a nonlinear Lipschitz (reconstruction) operator $\mathcal{R}$
\begin{align}
\label{Rspace}
\mathcal{R}: \quad \B_{R_{0}}^{[0,T]}(\P_n X^{\sigma}) \times \B_{R_{0}}^{[0,T]}(X^{\sigma}) \times \B_{R_{1}}^{[0,T]}(X^{\sigma+\veps}) \to \B_{R_{0}}^{[0,T]}(\Q_n X^{\sigma})
\end{align}
such that for any $U\in \B_{R_{0}}^{[0,T]}(X^{\sigma})$, $H_{1}\in \B_{R_{0}}^{[0,T]}(X^{\sigma})$ and $H_{2}\in \B_{R_{1}}^{[0,T]}(X^{\sigma+\veps})$ that satisfy
\begin{align}\label{UCabstract}
\left\{\begin{aligned}
\partial_t U&=AU+F(U+H_{1})+H_{2}, &&\textnormal{ on } [0,T], \\[3pt]
\bC U(t)&=0, &&\textnormal{ for }t\in [0,T],
\end{aligned}\right.
\end{align}
then $\Q_n U=\mathcal{R}(\P_{n}U,H_{1},H_{2})$.

Moreover, if~additionally $F$ satisfies \cref{assumFholom} then there exist $\eta,\eta_1>0$ such that, for any $U\in \B_{R_{0}}^{[0,T]}(X^{\sigma})$ that satisfies \eqref{UCabstract}, $\mathcal{R}$ extends holomorphically as
\begin{align*}
\mathcal{R}: \bigl( \P_nU+ \B_{\eta,\eta}^{[0,T]}(\P_n X^{\sigma})\bigr) \times \B_{R_{0},\eta}^{[0,T]}(X^{\sigma}) \times \B_{R_{1},R_{1}}^{[0,T]}(X^{\sigma+\veps}) \to \B_{R_{0},\eta_1}^{[0,T]}(\Q_n X^{\sigma}).
\end{align*}
\end{theorem}

The main idea is to consider the splitting
\begin{align*}
U=\P_n U+\Q_n U:=V+W,
\end{align*}
so the problem \eqref{UCabstract} translates into the following nonlinear observability system
\begin{align*}
\left\{\begin{array}{rl}
\partial_t V(t)&=AV(t)+\P_n F( V+W+H_{1})+\P_{n}H_{2}, \\[3pt]
\partial_t W(t)&=AW(t)+\Q_n F(V+W+H_{1})+\Q_{n}H_{2}, \\[3pt]
\bC V(t)&=-\bC W(t).
\end{array}\right.
\end{align*}
For a given bounded function $V\in C^0([0, T], \P_n X^\sigma)$, we~are interested in solving the nonlinear observability problem
\begin{align*}
\left\{\begin{array}{rl}
\partial_t W(t)&=AW(t)+\Q_n F(V+W+H_{1})+\Q_{n}H_{2}, \\[3pt]
\Pi_n\bC W &=- \Pi_n\bC V.
\end{array}\right.
\end{align*}
Note that the operator $\Pi_{n}$ is nonlocal in time in $[0,T]$.
\begin{proposition}\label{propabstract}
Let $R_{0}$ and $R_{1}>0$. Under Assumptions \ref{assumAA}, \ref{assumC} and \ref{assumF} (with $R_{0}$), there exists $n_{0}\in \N$ and $\eta>0$ such that for any $n\geq n_{0}$, for any $H_{1}\in \B_{5R_{0}/2}^{[0,T]}(X^{\sigma})$, $H_{2}\in \B_{R_{1}}^{[0,T]}(X^{\sigma+\veps})$ and $G\in \mathbb{B}_{\eta}(L^2([0, T], X^{\sigma}))$, there exists a unique solution $W\in \B_{R_{0}}^{[0,T]}(\Q_n X^{\sigma})$ to
\begin{align}\label{eqWG}
\left\{\begin{array}{rl}
\partial_t W&=AW+\Q_n F(W+H_{1})+\Q_{n}H_{2}, \\[3pt]
\Pi_n\bC W(t) &=\Pi_nG.
\end{array}\right.
\end{align}
This defines a nonlinear Lipschitz operator $\mathcal{\widetilde{R}}$
\begin{align}
\left\{\begin{array}{rcl}
\B_{5R_{0}/2}^{[0,T]}(X^{\sigma}) \!\times\! \B_{R_{1}}^{[0,T]}(X^{\sigma+\veps}) \!\times\! \mathbb{B}_{\eta}(L^{2}([0,T],X^{\sigma})) & \dpl\to& C^0([0, T], \Q_{n }X^{\sigma})\\[3pt]
(H_{1},H_{2},G)&\dpl\mto& W:=\mathcal{\widetilde{R}}(H_{1},H_{2},G).
\end{array}\right.
\end{align}
Moreover, if~additionally $F$ satisfies \cref{assumFholom}, then there exists $\eta>0$ such that~$\mathcal{\widetilde{R}}$ extends holomorphically as
\begin{align*}
\mathcal{\widetilde{R}}: \B_{5R_{0}/2,\eta}^{[0,T]}(X^{\sigma}) \times \B_{R_{1}, R_{1}}^{[0,T]}(X^{\sigma+\veps})\times \mathbb{B}_{\eta,\eta}(L^{2}([0,T],X^{\sigma})) \to C^0([0, T], \Q_{n }X^{\sigma}_{\mathbb{C}}).
\end{align*}
\end{proposition}
\begin{proof}[Proof of \cref{thmabstract} from \cref{propabstract}]
Let $\chi\in C^{\infty}_{0}(\R,[0,1])$ be supported in $[-1,1]$ such that $\chi(s)=1$ for $s\in [-1/2,1/2]$. We~define the operator $\mathcal{R}$ by the formula
\begin{align}\label{defR}
\mathcal{R}(V,H_{1},H_{2})=\mathcal{\widetilde{R}}(V+H_{1},H_{2},- \chi(\eta^{-1} \nor{\bC V}{L^{2}([0,T],X^{\sigma})}) \bC V),
\end{align}
where $\eta$ is the one in \cref{propabstract}. If $V$ is complex-valued, this is to be understood as $\chi(\eta^{-1} \nor{\bC V}{L^{2}([0,T],X^{\sigma}_{\mathbb{C}})})$. We~prove that it is well-defined and satisfies the requirements of \cref{thmabstract}.

We first observe that $\chi(\eta^{-1} \nor{\bC V}{L^{2}([0,T],X^{\sigma})}) =0$ if $\nor{ \bC V}{L^{2}([0,T],X^{\sigma})}\geq \eta$, so that we always have $\bnor{ \chi(\eta^{-1} \nor{ \bC V}{L^{2}([0,T],X^{\sigma})}) \bC V}{L^{2}([0,T],X^{\sigma})}\leq \eta$. In~particular, $\mathcal{R}$~is well-defined from the spaces made precise in \eqref{Rspace}.

We now~check that it satisfies the required properties. Let $U$ be as in the theorem and a solution of \eqref{UCabstract}. We~want to check that for $n$ large enough (only depending on the parameters), we~can impose $\Q_n U$ to be small in $C^{0}([0,T],X^{\sigma})$ and $\bC \P_n U$ to be small in $X^{\sigma}$. We~can apply \cref{prop:abs-prop-reg} to $U$ to propagate regularity, resulting in $\norm{U}_{C^0([0, T], X^{\sigma+\veps})}\leq R_{2}$ for some $R_{2}>0$. We~have
\begin{align}
\label{unifboundU}
\nor{ \Q_n U}{C^{0}([0,T],X^{\sigma})} &\leq \frac{1}{\inn{\lambda_n}^{\veps}} \nor{ \Q_n U}{C^{0}([0,T],X^{\sigma+\veps})}\leq \frac{R_{2}}{\inn{\lambda_n}^{\veps}} .
\end{align}
For such a solution satisfying \eqref{UCabstract}, we~have $\bC U=0$ and, in particular, $\bC \Q_n U=-\bC \P_n U$, so that
\begin{align*}
\nor{\bC \P_n U}{L^{2}([0,T],X^{\sigma})}=\nor{\bC \Q_n U}{L^{2}([0,T],X^{\sigma})}&\leq T^{1/2}\nor{\bC}{\mc{L}(X^\sigma)}\nor{ \Q_n U}{C^{0}([0,T],X^{\sigma})}\\
&\leq \frac{T^{1/2}\nor{\bC}{\mc{L}(X^\sigma)}R_{2}}{\inn{\lambda_n}^{\veps}} .
\end{align*}
Since $\lambda_n\to +\infty$ as $n\to +\infty$, we~can find $n\geq n_0$, so that the term above can be made smaller that $\eta/4$. In~particular, defining $W:= \mathcal{R}(\P_{n}U,H_{1},H_{2})$ and by definition of $\mathcal{R}$,
\begin{align*}
W= \mathcal{\widetilde{R}}(\P_{n}U+H_{1},H_{2},- \bC \P_{n}U)
\end{align*}
is the unique solution in $\B_{R_{0}}^{[0,T]}(\Q_n X^{\sigma})$ of
\begin{align*}
\left\{\begin{array}{rl}
\partial_t W&=AW+\Q_n F(W+\P_{n}U+H_{1})+\Q_{n}H_{2}, \\[3pt]
\Pi_n\bC W &=-\Pi_n \bC \P_{n}U.
\end{array}\right.
\end{align*}
Also, since $\bC \P_{n}U=- \bC \Q_{n}U $ we note that $\Q_n U$ is solution of
\begin{align*}
\left\{\begin{array}{rl}
\partial_t \Q_n U&=A\Q_n U+\Q_n F(\Q_n U+\P_n U+H_{1})+\Q_{n}H_{2}, \\[3pt]
\Pi_n\bC \Q_n U &=-\Pi_n \bC \P_{n}U.
\end{array}\right.
\end{align*}
In particular, since $ \nor{\bC \P_n U}{L^{2}([0,T],X^{\sigma})}\leq \eta/4$ and by the various bounds above, we~have
\[
\P_n U+H_{1}\in \B_{2R_{0}}^{[0,T]}(X^{\sigma}),\quad H_{2}\in \B_{R_{1}}^{[0,T]}(X^{\sigma+\veps})\quand-\bC \P_{n}U\in \mathbb{B}_{\eta}(L^2([0, T], X^{\sigma})).
\]
By definition of $\mathcal{\widetilde{R}}$, this implies $\Q_n U= \mathcal{\widetilde{R}}(\P_n U+H_{1},H_{2},- \bC \P_{n}U)$. So we conclude $\Q_{n}U=W=\mathcal{R}(\P_{n}U,H_{1},H_{2})$ as expected.

Concerning the holomorphic extension, since $\mathcal{\widetilde{R}}$ has a holomorphic extension, as proved in \cref{propabstract}, using formula \eqref{defR} and composition of function, it~is enough to prove that $V\mto \chi(\eta^{-1} \nor{\bC V}{L^{2}([0,T],X^{\sigma}_{\mathbb{C}})})$ is constant equal to $1$ and therefore holomorphic in some neighborhood $\P_nU+ \B_{\eta_{1},\eta_{1}}^{[0,T]}(\P_n X^{\sigma})$ that would be included in $\B_{3R_0/2,\eta/2}^{[0,T]}(X^{\sigma})$. This can be obtained for $\eta_{1}$ small enough since $ \nor{\bC \P_n U}{L^{2}([0,T],X^{\sigma})}\leq \eta/4$. This gives the result up to renaming $\eta_{1}$ by $\eta$.
\end{proof}

\begin{proof}[Proof of \cref{propabstract}]
Using \cref{lmCauchyobs} (recall that $\FL$ denotes a linear flow map), we~are looking for $W$ solution of\vspace*{-3pt}
\begin{align*}
W=\FL (G,F(W+H_{1})+H_{2}).
\end{align*}
So, it~is natural to define the nonlinear operator $\Phi_{n,G,H_{1},H_{2}}$ defined by\vspace*{-3pt}
\begin{align}\label{defPhi}
\Phi_{n,G,H_{1},H_{2}}(W):=\FL (G,F(W+H_{1})+H_{2})
\end{align}
which we will later prove to be well-defined from suitable balls of $C^0([0, T], \Q_n X^\sigma)$ to $C^0([0, T], \Q_n X^\sigma)$. To keep notations reasonable, we~will write $\Phi=\Phi_{n,V,G,H_{1},H_{2}}$ keeping in mind all the dependence. The goal will be to find a fixed point of the operator $\Phi$ in a small ball that will also satisfy \eqref{eqWG}.

Following Hale-Raugel \cite[Th.\,2.14]{HR03} and Joly-Laurent \cite[Th.\,10.1]{2020:joly-laurent:decay-nlw-no-gcc}, we~divide the proof in three steps.

\subsubsection*{Step 1. High-frequency fixed point: real case} For this part of the proof, we~consider all the functions involved to be real-valued, in~the sense that we consider the real vector space $X^{\sigma}$. Recall that throughout, we~work with the closed balls $\mathbb{B}$ and $\B$ defined at the beginning of Section~\ref{s:abstrIntro}.

We fix $H_{1}\in \B_{5R_{0}/2}^{[0,T]}(X^{\sigma})$, $H_{2}\in \B_{R_{0}}^{[0,T]}(X^{\sigma+\veps})$ and $G\in \mathbb{B}_{\eta}(L^2([0, T], X^{\sigma}))$, but we will make estimates uniform when these functions are in these sets. We~want to find a fixed point in $W\in \B_{R_{0}}^{[0,T]}(\Q_n X^{\sigma})$ by the fixed point theorem of Banach. We~prove that $\Phi$ is a contraction in this set.

Let $W\in \B_{R_{0}}^{[0,T]}(\Q_n X^{\sigma})$. Estimate \eqref{estimFL} of \cref{lmCauchyobs} gives for some constants uniform in $n\in \N$\vspace*{-5pt}
\begin{multline*}
\nor{\Phi(W)}{C^{0}([0,T],\Q_{n}X^{\sigma})}=\nor{\FL(G,F(W+H_{1})+H_{2})}{C^{0}([0,T],\Q_{n}X^{\sigma})}\\
\leq C \nor{\Pi_{n}G}{L^{2}([0,T],X^{\sigma})}+C\nor{\Q_{n}\left[ F(W+H_{1})+H_{2}\right]}{L^{1}([0,T],X^{\sigma})}.
\end{multline*}
Since $\Pi_{n}$ is a projection on $L^2([0,T],X^{\sigma})$, by~assumption we have\vspace*{-3pt}
\[
\nor{\Pi_{n}G}{L^{2}([0,T],X^{\sigma})}\leq \nor{G}{L^{2}([0,T],X^{\sigma})}\leq \eta.
\]
To estimate the second term, since $\nor{W+H_{1}}{C^{0}([0,T],X^{\sigma})}\leq 3R_{0}$, we~can use \eqref{boundF} and obtain\vspace*{-3pt}
\begin{equation}\label{estimFinPhi}
\begin{aligned}
\nor{\Q_{n}\left[ F(W+H_{1})\!+\!H_{2}\right]}{L^{1}([0,T],X^{\sigma})}&\leq \frac{1}{\inn{\lambda_n}^{\veps}} \nor{F(W+H_{1})\!+\!H_{2}}{L^{1}([0,T],X^{\sigma+\veps})}\\
&\leq \frac{T}{\inn{\lambda_n}^{\veps}}(C+R_{1}).
\end{aligned}
\end{equation}
Summing up, the previous estimates gives
\begin{align*}
\nor{\Phi(W)}{C^{0}([0,T],\Q_{n}X^{\sigma})}\leq C\eta+ \frac{1}{\inn{\lambda_n}^{\veps}} \left(C+R_{1}\right).
\end{align*}
Concerning the difference, similar estimates give for $W, W'\in \B_{R_{0}}^{[0,T]}(\Q_n X^{\sigma})$, using instead the second estimate of \eqref{boundF}
\begin{align}
\nonumber\nor{\Phi(W)-\Phi(W')}{C^{0}([0,T],\Q_{n}X^{\sigma})}&=\nor{\FL(0,F(W+H_{1})-F(W'+H_{1}))}{C^{0}([0,T],\Q_{n}X^{\sigma})}\\
\nonumber&\leq C\nor{\Q_{n}\left[ F(W+H_{1})-F(W'+H_{1}))\right]}{L^{1}([0,T],X^{\sigma})}\\[-7pt]
\label{bounddiffW}\\[-10pt]
\nonumber&\leq \frac{C}{\inn{\lambda_n}^{\veps}} \nor{F(W+H_{1})-F(W'+H_{1}))}{L^{1}([0,T],X^{\sigma+\veps})}\\
\nonumber&\leq \frac{CCT}{\inn{\lambda_n}^{\veps}} \nor{W-W'}{C^{0}([0,T],X^{\sigma})}.
\end{align}
In particular, if~$\eta$ is chosen small enough and $n_{0}$ is large enough, then $\Phi$ reproduces the closed ball $\B_{R_{0}}^{[0,T]}(\Q_n X^{\sigma})$ and is contracting.

\subsubsection*{Step 2. Extension to the complex case} The proof is very similar to the real case. We~define $\Phi$ with the same formula as \eqref{defPhi}, but instead consider $W\in \B_{R_{0},\eta_{1}}^{[0,T]}(\Q_n X^{\sigma})$, $H_{1}\in \B_{5R_{0}/2,\eta}^{[0,T]}(X^{\sigma})$, $H_{2}\in \B_{R_{1}, R_{1}}^{[0,T]}(X^{\sigma+\veps})$ and $G\in \mathbb{B}_{\eta,\eta}(L^{2}([0,T],X^{\sigma}))$, where $\eta_{1}>0$ small is to be determined later.

The flow map $\FL$ is linear, so the extension to complex-valued vectors is clear, see \cite[Th.\,3]{BS:71-polynomials}. So, we~only need to check that the valuation $F(W+H_{1})$ makes sense and the contraction in $\B_{R_{0},\eta_{1}}^{[0,T]}(\Q_n X^{\sigma})$ of $\Phi$ is still true with these parameters, up to making $\eta$ small and $n$ large.

If $W\in \B_{R_{0},\eta_{1}}^{[0,T]}(\Q_n X^{\sigma}), H_{1}\in \B_{5R_{0}/2,\eta}^{[0,T]}(X^{\sigma})$, then
\[
W+H_{1}\in \B_{7R_{0}/2,\eta_{1}+\eta}^{[0,T]}(X^{\sigma})\subset \B_{4R_{0},\delta}^{[0,T]}(X^{\sigma})
\]
if we have $\eta_{1}+\eta\leq \delta$, so that $F(W+H_{1})$ is well-defined and satisfies similar estimates as in the real case, thanks to \eqref{boundFanalytic} in \cref{assumFholom}. It~remains to check $\Phi(W)\in \B_{R_{0},\eta}^{[0,T]}(\Q_n X^{\sigma})$. We~still have
\begin{align*}
\nor{\Pi_{n}G}{L^{2}([0,T],X_\mathbb{C}^{\sigma})}\leq 2\eta,
\end{align*}
when $G\in \mathbb{B}_{\eta,\eta}(L^{2}([0,T],X^{\sigma}))$, while a similar estimate as in \eqref{estimFinPhi} gives
\begin{align*}
\nor{\Q_{n}\left[ F(W+H_{1})+H_{2}\right]}{L^{1}([0,T],X_\mathbb{C}^{\sigma})}&\leq \frac{1}{\inn{\lambda_n}^{\veps}} \left(C+2R_{1}\right).
\end{align*}
Finally, the linear estimate \eqref{estimFL} of \cref{lmCauchyobs} still holds in the complex-valued case, so we~get
\begin{align*}
\nor{\Phi(W)}{C^{0}([0,T],\Q_{n}X_\mathbb{C}^{\sigma})}\leq 2\eta+ \frac{1}{\inn{\lambda_n}^{\veps}} \left(C+2R_{1}\right).
\end{align*}
The bound \eqref{bounddiffW} holds in the complex case without modification. So, we~obtained that if $\eta_{1}+\eta\leq \delta$, $4\eta\leq \eta_{1}\leq R_{0}$ and $n$ is large enough, $\Phi$ is a contraction on $\B_{R_{0},\eta_{1}}^{[0,T]}(\Q_n X^{\sigma})$. This gives a unique fixed point in the complex case.

\subsubsection*{Step 3. Regularity of the fixed point} We prove that for fixed $n$, the map
\[
\widehat{\Phi}: (W, G, H_1, H_2)\mto \Phi_{n, G, H_1, H_2}(W)
\]
is Lispchitz of its arguments under \cref{assumF} and holomorphic when \cref{assumFholom} is made.
Recalling Definition \eqref{defPhi} and that $\FL$ is linear in its arguments, we~see that, by~composition, it~will be enough to verify that $(W,H_1)\mto F(W+H_{1})$ is Lipschitz or holomorphic, depending on the assumption made.

First of all, observe that the same argument to get estimates \eqref{bounddiffW}, shows that $(W,H_1)\mto F(W+H_{1})$ is Lipschitz under \cref{assumF}, and by linearity of the flow map $\mc{F}_n$, so is the map
\begin{multline*}
\widehat{\Phi}: \B_{R_0}^{[0, T]}(\Q_n X^{\sigma})\times \mathbb{B}_{\eta}(L^{2}([0,T],X^{\sigma}))\times \B_{2R_{0}}^{[0,T]}(X^{\sigma}) \times \B_{R_{1}}^{[0,T]}(X^{\sigma+\veps})\\
\to C^0([0, T], \Q_{n }X^{\sigma}).
\end{multline*}
Now, under \cref{assumFholom}, the same argument as before shows that the complex extension, denoted by the same letter $\widehat{\Phi}$,
\begin{multline*}
\widehat{\Phi}: \B_{R_0, \eta}^{[0, T]}(\Q_n X^{\sigma})\times \mathbb{B}_{\eta,\eta}(L^{2}([0,T],X^{\sigma}))\times \B_{2R_{0},\eta}^{[0,T]}(X^{\sigma}) \times \B_{R_{1}, R_{1}}^{[0,T]}(X^{\sigma+\veps})\\ \to C^0([0, T], \Q_{n }X_\mathbb{C}^{\sigma})
\end{multline*}
is Lipschitz, hence continuous. In~view of \cite[Prop.\,8.10]{Muj86} and that $\widehat{\Phi}$ is linear in the variables $G$ and $H_2$ separately, the main task is to show that the map
\begin{align}
\left\{\begin{array}{rcl}
\B_{4R_{0}, \eta}^{[0,T]}(X_{\mathbb{C}}^{\sigma}) \times \B_{2R_{0}, \eta}^{[0,T]}(X^{\sigma}) & \dpl\to& C^0([0, T], X_\mathbb{C}^{\sigma+\veps})\\[3pt]
(W, H_1 )&\dpl\mto& F(W+H_1)
\end{array}\right.
\end{align}
is holomorphic on each variable separately. Since $F$ is holomorphic when defined from some balls to $X^\sigma$ to $X^{\sigma+\veps}$, the only point to check is that it implies the same result as a map on time-dependent functions in $C^0([0,T],X^\sigma)$.

Fix $W\in \text{Int}(\B_{R_0, \eta}^{[0, T]}(\Q_nX^\sigma))$ and $H_1\in \text{Int}(\B_{5R_0/2, \eta}^{[0, T]}(X^\sigma))$, and $\ell>0$ such that $B_{C^0([0,T],X_{\C}^\sigma)}(W+H_1,\ell)\subset \B_{4R_0, 2\eta}^{[0, T]}(X^\sigma))$. Let $C$ be the bound on $F$ in \cref{assumFholom}. For all
\begin{itemize}
\item
$K_W\in C^0([0, T], \Q_n X_\mathbb{C}^\sigma)$,
\item
$K_H\in C^0([0, T], X_\mathbb{C}^\sigma)$ with $\norm{K_W+K_H}_{C^0([0, T], X_\mathbb{C}^\sigma)}\leq \ell/4$,
\item
and $t\in [0,T]$,
\end{itemize}
we~have $(W+H_1+K_W+K_H)(t) \in \mathbb{B}_{4R_{0},2\delta}(X^{\sigma})$ (recall $\eta<\delta$). Being $F$ holomorphic, by~applying a Taylor expansion (see \cite[Th.\,7.1, Cor.\,7.3]{Muj86}) together with Cauchy estimates (see \cite[Cor.\,7.4]{Muj86}), we~have the bound, uniform in $t\in [0,T]$,
\begin{align*}
\lVert F(W(t)+H_1(t)+K_W(t)+K_H(t))&-F(W(t)+H_1(t))\\
&-\delta F(W(t)+H_1(t); K_W(t)+K_H(t)) \rVert_{ X_\mathbb{C}^{\sigma+\veps}}\\
&\leq \dfrac{4C\norm{K_W(t)+K_H(t)}_{X_\mathbb{C}^\sigma}^2}{\ell(\ell-2\norm{K_W(t)+K_H(t)}_{X_\mathbb{C}^\sigma})} \\
&\leq \dfrac{4C\norm{K_W+K_H}_{C^0([0, T], X_\mathbb{C}^\sigma)}^2}{\ell(\ell-2\norm{K_W+K_H}_{C^0([0, T], X_\mathbb{C}^\sigma)})}\\
&\leq \dfrac{8C\norm{K_W+K_H}_{C^0([0, T], X_\mathbb{C}^\sigma)}^2}{\ell^2},
\end{align*}
where
\begin{align*}
\delta F(Z; K):=\lim_{s\to 0}\dfrac{1}{s}(F(Z+sK)-F(Z)).
\end{align*}
Observe that the differential of $F$ can be easily extended from $C^0([0, T], X_\mathbb{C}^\sigma)$ into $C^0([0, T], X_\mathbb{C}^{\sigma+\veps})$. Since, the previous estimate is uniform in $t\in [0,T]$, we~can write
\begin{multline}
\bigl\Vert F(W+H_1+K_W+K_H)-F(W+H_1)\\
\shoveright{-\delta F(W+H_1; K_W+K_H)\bigr\Vert_{C^0([0, T], X_\mathbb{C}^{\sigma+\veps})}}\\
\leq \dfrac{8C\norm{K_W+K_H}_{C^0([0, T], X_\mathbb{C}^\sigma)}^2}{\ell^2}.\label{prop:proof:frechet-1}
\end{multline}

Actually, the previous estimate \eqref{prop:proof:frechet-1} holds uniformly in a ball around $(W, H_1)$ of radius $\ell/4$. Consequently, we~establish that the map
\begin{align*}
\left\{\begin{array}{rcl}
\B_{R_{0}, \eta}^{[0,T]}(X^{\sigma}) \times \B_{5/2R_{0}, \eta}^{[0,T]}(X^{\sigma}) &\dpl\to& C^0([0, T], X_\mathbb{C}^{\sigma+\veps})\\[3pt]
(W, H_1 )&\dpl\mto& F(W+H_1)
\end{array}\right.
\end{align*}
is Fréchet differentiable (in the complex sense), hence holomorphic as an application of Theorem \cite[Th.\,13.16]{Muj86}. The estimates in \cref{lmCauchyobs} show that the map $\FL$ is bilinear continuous with respect to each of its parameters and therefore has an obvious holomorphic extension. Since the embedding $\iota: C^0([0,T],X_{\mathbb{C}}^{\sigma+\veps})\to L^1([0,T],X_{\mathbb{C}}^{\sigma})$ is linear, an application of the chain rule (see \cite[Th.\,13.6]{Muj86}) shows that the map $\widehat{\Phi}$ is holomorphic since it is defined by $\Phi_{n,G,H_{1},H_{2}}(W):=\FL (G,\iota\left(F(W+H_{1})+H_{2}\right)) $. The conclusion that the nonlinear reconstruction operator $\mathcal{\widetilde{R}}$ is holomorphic follows as a direct application of \cite[Th.\,2.2]{1982:chow-hale:bifurcation-theory} of regularity of fixed points with respect to parameters.
\end{proof}

\subsection{Analyticity in time of the observed solution} Here we prove the main theorem of this section, that is the abstract \cref{thmabstractanalyticintro} of analytic regularity stated in the introduction.
\begin{proof}[Proof of \cref{thmabstractanalyticintro}]
First, since $y\mto \max_{t\in [0,T^{*}]}\nor{\Im H_{1}(t+iy)}{X^{\sigma}}$ is a continuous function on $[-\mu,\mu]$ that is equal to zero at $y=0$, there exists $0<\mu'<\mu$ such that $\max_{t\in [0,T]}\left|\Im H_{1}(t+iy)\right| \leq \eta$ for $y\in [-\mu',\mu']$, where $\eta$ is given by \cref{thmabstract}. We~can do the same for $H_{2}$. We~denote $R_{1}=\max_{z\in [0,T^{*}]+i[-\mu,\mu]}\nor{H_{2}(z)}{X^{\sigma+\veps}}$.

Let $n\in \N$, $\eta>0$ and $\mathcal{R}$ be given by \cref{thmabstract}, with the chosen $R_0$ and $R_{1}$. Let~$\nu$ be such that $0<\nu<T^*-T$.

With these two choices, if~we define for $s\in [-\nu, T^*-T-\nu]$, $H_{1}^{s}$ as the map $t\mto H_{1}(t+\nu+s)$, we~have $H_{1}^{s} \in \B_{R_{0},\eta}^{[0,T]}(X^{\sigma})$ for any $s\in [-\nu, T^*-T-\nu]$. Moreover, the map $s\mto H_1^s$ is continuous. Indeed, since $t\in [0, T^*]\mto H_1(t)\in X^\sigma$ is a continuous map defined on a compact set, it~is uniformly continuous. Hence, for every $\veps>0$ and any $t\in [0, T]$, there exists $\delta>0$ such that for any $s$, $s_0\in [-\nu, T^*-T-\nu]$ with $|(s+t)-(s_0+t)|\leq \delta$, then $\norm{H_1(t+\nu+s)-H_1(t+\nu+s_0)}_{X^\sigma}\leq \veps$. Since the later property does not depend on $t$, we~can take $\sup_{t\in [0, T]}$ in the last inequality to obtain that $\norm{H_1^s-H_1^{s_0}}_{C^0([0, T], X^\sigma)}\leq \veps$. We~do the same for $H_2$, so we have $H_2^s\in \B_{R_{1},R_{1}}^{[0,T]}(X^{\sigma+\veps})$ for any $s\in [-\nu, T^*-T-\nu]$ and moreover $s\mto H_2^s$ is continuous.

For $U$, we~can define similarly $U^s\in \B_{R_{0},\eta}^{[0,T]}(X^{\sigma})$ for any $s\in [-\nu, T^*-T-\nu]$. We~can also decompose $U=\P_n U+\Q_n U=:V+W$ and $U^s=V^s+W^s$. For any fixed $s\in [-\nu, T^*-T-\nu]$, the assumption \eqref{UCabstractT*intro} implies
\begin{align}\label{UCabstractwiths}
\left\{\begin{aligned}
\partial_t U^s&=AU^s+F(U^s+H_{1}^s)+H_{2}^s&& \textnormal{on } [0,T],\\
\bC U^s(t)&=0 &&\textnormal{for }t\in [0,T].
\end{aligned}\right.
\end{align}
In particular, for any fixed $s\in [-\nu, T^*-T-\nu]$, \cref{thmabstract} gives
\begin{align}
\label{equalWs}
W^s=\Q_n U^s=\mathcal{R}(\P_{n}U^s,H_{1}^s,H_{2}^s)=\mathcal{R}(V^s,H_{1}^s,H_{2}^s),
\end{align} the equality being meant in $C^0([0,T], \Q_n X^{\sigma})$.

But if we denote
\[
G=\P_n F(U+H_{1})+\P_n H_{2}=\P_n F(V+W+H_{1})+\P_n H_{2}\in C^0([0,T^*],\P_n X^{\sigma}),
\]
then $V\in C^0([0,T^*],\P_n X^{\sigma})$ is solution of
\begin{align*}
\partial_t V=\P_n AV+G& \textnormal{ on } [0,T^{*}].
\end{align*}
With the related notations, \cref{lmtranslat} implies
\begin{align}
\label{eqVG}
\partial_s V^s=\P_n AV^s+G^s& \textnormal{ on } [-\nu, T^*-T-\nu].
\end{align}
But for any fixed $s\in [-\nu, T^*-T-\nu]$, we~have, using \eqref{equalWs},
\begin{align*}
G^s&=\P_n F(V^s+W^s+H_{1}^s)+\P_n H_{2}^s\\ &=\P_n F(V^s+\mathcal{R}(V^s,H_{1}^s,H_{2}^s)+H_{1}^s)+\P_n H_{2}^s
\end{align*}
with equality in $C^0([0,T], \P_n X^{\sigma})$.

It is then natural to consider the map
\begin{align*}
\Tilde{F}: [-\nu, T^*-T-\nu]\times \B_{R_{0}}^{[0,T]}(\P_n X^{\sigma})\to C^0([0, T], \P_n X^\sigma)
\end{align*}
defined by
\begin{align*}
\Tilde{F}(s,\Tilde{V})=\P_n F(\Tilde{V}+\mathcal{R}(\Tilde{V},H_{1}^s,H_{2}^s)+H_{1}^s)+\P_n H_{2}^s.
\end{align*}
This map is well-defined, continuous and locally Lipschitz (with respect to the second variable). We~observe that $H_1^s\in \B_{R_0}^{[0, T]}(X^\sigma)$, $H_2^s\in \B_{R_{1}}^{[0, T]}(X^{\sigma+\veps})$ for all \hbox{$s\in [-\nu, T^*-T-\nu]$} and $\Tilde{V}\in \B_{R_{0}}^{[0,T]}(\P_n X^{\sigma})$ implies, by~construction, that
\[
\Tilde{V}+\mathcal{R}(\Tilde{V},H_{1}^s,H_{2}^s)+H_{1}^s\in \B_{3R_0}^{[0, T]}(X^\sigma).
\]
Therefore $F(\Tilde{V}+\mathcal{R}(\Tilde{V},H_{1}^s,H_{2}^s)+H_{1}^s)$ is well-defined and so is $\Tilde{F}$. Recall that $\mc{R}$ is Lipschitz on its variables, that is, there exists $L>0$ such that
\begin{multline*}
\norm{\mc{R}(\Tilde{V}, H_{1}^{s}, H_2^{s})-\mc{R}(\Tilde{V}^*, H_{1}^{s^*}, H_2^{s^*})}_{C^0([0, T], X^\sigma)}\leq L\big( \norm{\Tilde{V}-\Tilde{V}^*}_{C^0([0, T], X^\sigma)}\\+\norm{H_1^{s}-H_1^{s^*}}_{C^0([0, T], X^\sigma)}+\norm{H_2^{s}-H_2^{s^*}}_{C^0([0, T], X^{\sigma+\veps})} \big),
\end{multline*}
for all $\Tilde{V}$, $\Tilde{V}^*\in \B_{R_{0}}^{[0,T]}(\P_n X^{\sigma})$ and $s$, $s^*\in [-\nu, T^*-T-\nu]$. The continuity of $\Tilde{F}$ then follows by the continuity of $s\mto (H_1^s, H_2^s)$, \cref{assumFholom} and by algebra of continuous maps. The same inequality along with \cref{assumFholom} shows that
\begin{align*}
\norm{\Tilde{F}(s, \Tilde{V})-\Tilde{F}(s, \Tilde{V}^*)}_{C^0([0, T], \P_n X^\sigma)}\leq C\norm{\Tilde{V}-\Tilde{V}^0}_{C^0([0, T], \P_n X^\sigma)},
\end{align*}
for any $\Tilde{V}$, $\Tilde{V}^*\in \B_{R_{0}}^{[0,T]}(\P_n X^{\sigma})$ and $s\in [-\nu, T^*-T-\nu]$.

Moreover, we~claim that for any $s_0\in (-\nu,T^*-T-\nu)$, there exists $\rho>0$ such that it admits a holomorphic extension (in both variables)
\begin{align*}
\Tilde{F}: B_{\mathbb{C}}(s_0, \rho)\times \bigl[V^{s_0}+\B_{\eta,\eta}^{[0,T]}(\P_n X^{\sigma})\bigr]\to C^0([0, T], \P_n X_\mathbb{C}^\sigma)
\end{align*}
given by
\begin{align*}
\Tilde{F}(z,\Tilde{V})=\P_n F(\Tilde{V}+\mathcal{R}(\Tilde{V},H_{1}^z,H_{2}^z)+H_{1}^z)+\P_n H_{2}^z.
\end{align*}
By \cref{thmabstract}, around $V^{s_0}=\P_n U^{s_0}$, the map $\mc{R}$ extends holomorphically as
\begin{align*}
\mathcal{R}: \bigl[ V^{s_0}+ \B_{\eta,\eta}^{[0,T]}(\P_n X^{\sigma})\bigr] \times \B_{R_{0},\eta}^{[0,T]}(X^{\sigma}) \times \B_{R_{1},R_{1}}^{[0,T]}(X^{\sigma+\veps}) \to C^0([0, T], \Q_{n }X_\mathbb{C}^{\sigma}).
\end{align*}
We need to argue that, the map $z\mto H_1^z$ is holomorphic from $B_{\mathbb{C}}(s_0, \rho)$ to $C^0([0, T], X_\mathbb{C}^\sigma)$ for $\rho>0$ small enough, and that the same holds for $z\mto H_2^z$ from $B_{\mathbb{C}}(s_0, \rho)$ with values in $C^0([0, T], X_\mathbb{C}^{\sigma+\veps})$. Indeed, since
\begin{align*}
z\in (0, T^*)+i(-\mu', \mu')\mto H_1(z)\in X_\mathbb{C}^\sigma
\end{align*}
is holomorphic, for each $s_0\in [-\nu, T^*-T-\nu]$ we can find $\rho>0$ such that, for any $t\in [0, T]$, the map
\begin{align*}
z\in B_\mathbb{C}(s_0, \rho)\mto H_1(z+t+\nu)\in X_\mathbb{C}^\sigma
\end{align*}
is holomorphic. As we did in the proof of \cref{propabstract}, by~using Cauchy estimates and the bound on $H_1(z)$, for each $t\in [0, T]$ and any $z_0\in B_\mathbb{C}(s_0, \rho)$, we~can find $\ell>0$ such that
\begin{multline*}
\norm{H_1(z+h+t+\nu)-H_1(z+t+\nu)-\delta H_1(z+t+\nu, h)}_{X_\mathbb{C}^\sigma}\\
\leq \dfrac{4(R_0+\eta)|h|^2}{\ell(\ell-2|h|)}\leq \dfrac{8(R_0+\eta)|h|^2}{\ell^2}
\end{multline*}
holds uniformly for all $z\in B_\mathbb{C}(z_0, \ell/4)$ and $B_\mathbb{C}(0, \ell/4)$. Moreover, the last estimate is independent of $t\in [0, T]$, so we actually have
\begin{align*}
\norm{H_1^{z+h}-H_1^z-\delta H_1(z+\cdot+\nu, h)}_{C^0([0, T], X_\mathbb{C}^\sigma)}\leq \dfrac{8(R_0+\eta)|h|^2}{\ell^2},
\end{align*}
uniformly for all $z\in B_\mathbb{C}(z_0, \ell/4)$ and $|h|\leq \ell/4$. The previous estimate shows that the map
\begin{align*}
z\in B_\mathbb{C}(s_0, \rho)\mto H_1^z\in \B_{R_0, \eta}^{[0, T]}(X^\sigma)
\end{align*}
is holomorphic. The argument works similarly to show that $z\mto H_2^z$ is holomorphic.

By composition of holomorphic functions and \cite[Prop.\,8.10]{Muj86} we get that
\begin{align*}
\left\{\begin{array}{rcl}
B_{\mathbb{C}}(s_0, \rho)\times \bigl[V^{s_0}+\B_{\eta,\eta}^{[0,T]}(\P_n X^{\sigma})\bigr]&\dpl\to& C^0([0, T], \Q_n X_\mathbb{C}^\sigma)\\[3pt]
(z,\Tilde{V})&\dpl\mto& \mathcal{R}(\Tilde{V},H_{1}^z,H_{2}^z)
\end{array}\right.
\end{align*}
is a well-defined and holomorphic map. That the extension $\Tilde{F}$ is holomorphic, follows from \cref{assumFholom}, algebra of holomorphic maps, and that $\P_n$ is a linear bounded operator.

Since we have proved that for any $s\in [-\nu, T^*-T-\nu]$, $G^s=\Tilde{F}(s,V^s)$, we~obtain that the equation \eqref{eqVG} verified by $V^s$ can be written as
\begin{align*}
\partial_s V^s=\P_n AV^s+\Tilde{F}(s,V^s)& \textnormal{ on } [-\nu, T^*-T-\nu].
\end{align*}
We then consider the following ODE, defined in the Banach space $C^0([0, T], \P_n X_\mathbb{C}^\sigma)$,
\begin{align}\label{thm:proof:eq:low-freq-evo-nlw}
\left\{\begin{array}{rl}
\partial_s \xi(s)&=\P_nA\xi(s)+\Tilde{F}(s, \xi(s)), \\[3pt]
\xi(s_0)&=V^{s_0}.
\end{array}\right.
\end{align}
Note that $C^0([0,T], \P_n X^\sigma)$ is of infinite dimension, but we can check that $\P_n A$ is a linear bounded operator on it. \cref{lmDuhamelCauchy} shows that $V^s$, which is a mild solution of \eqref{thm:proof:eq:low-freq-evo-nlw} (\ie satisfying \eqref{appendix:ode:ode-banach-1-duhamel}) is also a classical solution in the usual sense of ODE in Banach space.

Note that $\P_n A+\Tilde{F}$ is holomorphic from $ B_{\mathbb{C}}(s_0, \rho)\times \left[V^{s_0}+ \B_{\eta,\eta}^{[0,T]}(\P_n X^{\sigma})\right]$ into $C^0([0, T], \P_n X_{\mathbb{C}}^\sigma)$. Therefore, by~the classical theory of ODEs in Banach spaces \cite[Th.\,10.4.5]{1969:dieudonne:modern-analysis}, we~obtain a unique classical solution of \eqref{thm:proof:eq:low-freq-evo-nlw}
\begin{align*}
\xi:B_{\mathbb{C}}(s_0, \rho')\mto \Tilde{V}(z)\in C^0([0,T], \P_n X_{\mathbb{C}}^\sigma)
\end{align*}
for some $\rho'\in (0, \rho]$. Moreover, such a solution map inherits the regularity of the right-hand side of the ODE and so it is a holomorphic map. By~uniqueness of the solutions, we~have $\xi(s)=V^s$ for $s\in (s_0-\rho',s_0+\rho')$.

Moreover, the maps $z\mto \xi(z)$ and
\begin{align*}
\left\{\begin{array}{rcl}
B_{\mathbb{C}}(s_0, \rho)\times \bigl[V^{s_0}+ \B_{\eta,\eta}^{[0,T]}(\P_n X^{\sigma})\bigr]&\dpl\to& C^0([0, T], \P_n X_{\mathbb{C}}^\sigma)\\[3pt]
(z,\Tilde{V})&\dpl\mto&\Tilde{V}+ \mathcal{R}(\Tilde{V},H_{1}^z,H_{2}^z)
\end{array}\right.
\end{align*}
are both holomorphic in their respective spaces, and so is its composition\vspace*{-3pt}
\[
z\mto \xi(z)+\mathcal{R}(\xi(z),H_{1}^z,H_{2}^z).
\]

But for $s\in (s_0-\rho',s_0+\rho')$, we~have $\xi(s)=V^s$, so that
\[
\xi(s)+\mathcal{R}(\xi(s),H_{1}^s,H_{2}^s)=V^s+\mathcal{R}(V^s,H_{1}^s,H_{2}^s)=V^s+W^s=U^s,
\]
where we have used \eqref{equalWs} and the inclusion $(s_0-\rho',s_0+\rho')\subset (-\nu,T^*-T-\nu)$ for~$\rho'$ small enough. In~particular, the map $s\in (s_0-\rho',s_0+\rho')\mto U^s\in C^0([0, T], X^\sigma)$ is the restriction to a real interval of a holomorphic map, and hence is real analytic.

Since for any $t_0\in [0,T]$, the trace map $C^0([0, T], X^\sigma)\to X^\sigma$ is linear continuous, we~obtain by composition that the map\vspace*{-3pt}
\begin{align*}
\left\{\begin{array}{rcl}
(s_0-\rho',s_0+\rho')&\dpl\to& X^\sigma\\[3pt]
s&\dpl\mto&U^s(t_0)=U(t_0+\nu+s)
\end{array}\right.
\end{align*}
is real analytic. Recall that $\rho'$ depends on all the other parameters, but $\nu$ is an arbitrary number with $0<\nu<T^*-T$, $s_0$ is any number with $s_0\in (-\nu,T^*-T-\nu)$, while $t_0$ is arbitrary in $[0,T]$. It~means that $t\mto U(t)$, (which is well-defined for $t\in [0,T^*]$), is real analytic in a neighborhood of any $t_1$ of the form $t_1=t_0+\nu+s_0$ with $t_0$, $\nu$ and $s_0$ as before. It~is not hard to see that this implies that $t\mto U(t)$ is real analytic on $(0,T^*)$, as expected.
\end{proof}
